%----------------------------------------------------------------
%
%  File    :  thesis.tex
%
%  Authors :  Keith Andrews, IICM, TU Graz, Austria
%             Manuel Koschuch, FH Campus Wien, Austria
%			  Sebastian Ukleja, FH Campus Wien, Austria
%
%  Created :  22 Feb 96
%
%  Changed :  14 Oct 2020
%
%  For suggestions and remarks write to: sebastian.ukleja@fh-campuswien.ac.at
%
%----------------------------------------------------------------

% --- Setup for the document ------------------------------------

%Class for a book like style:
\documentclass[11pt,a4paper,oneside]{scrbook}
%For a more paper like style use this class instead:
%\documentclass[11pt,a4paper,oneside]{thesis}

%input encoding for windows in utf-8 needed for Ä,Ö,Ü etc..:
\usepackage[utf8]{inputenc}
%input encoding for linux:
%\usepackage[latin1]{inputenc}
%input encoding for mac:
%\usepackage[applemac]{inputenc}

\usepackage[english]{babel}
% for german use this line instead:
%\usepackage[ngerman]{babel}

%needed for font encoding
\usepackage[T1]{fontenc}

% want Arial? uncomment next two lines...
%\usepackage{uarial}
%\renewcommand{\familydefault}{\sfdefault}

%some formatting packages
\usepackage[bf,sf]{subfigure}
\renewcommand{\subfigtopskip}{0mm}
\renewcommand{\subfigcapmargin}{0mm}

%For better font resolution in pdf files
\usepackage{lmodern}

\usepackage{url}

%\usepackage{latexsym}

\usepackage{geometry} % define pagesize in more detail


\usepackage{colortbl} % define colored backgrounds for tables

\usepackage{courier} %for listings
\usepackage{listings} % nicer code formatting
\lstset{basicstyle=\ttfamily,breaklines=true}

\usepackage{graphicx}
  \pdfcompresslevel=9
  \pdfpageheight=297mm
  \pdfpagewidth=210mm
  \usepackage[         % hyperref should be last package loaded
    pdftex, 		   % needed for pdf compiling, DO NOT compile with LaTeX
    bookmarks,
    bookmarksnumbered,
    linktocpage,
    pagebackref,
    pdfview={Fit},
    pdfstartview={Fit},
    pdfpagemode=UseOutlines,                 % open bookmarks in Acrobat
  ]{hyperref}
\DeclareGraphicsExtensions{.pdf,.jpg,.png}
\usepackage{bookmark}

\usepackage[title]{appendix}

%paper format
\geometry{a4paper,left=30mm,right=25mm, top=30mm, bottom=30mm}

% --- Settings for header and footer ---------------------------------
\usepackage{scrlayer-scrpage}
\clearscrheadfoot
\pagestyle{scrheadings}
\automark{chapter}

%Left header shows chapter and chapter name, will not display on first chapter page use \ihead*{\leftmark} to show on every page
\ihead{\leftmark}
%\ohead*{\rightmark}	%optional right header
\ifoot*{Student}		%left footer shows student name
\ofoot*{\thepage}		%right footer shows pagination
%---------------------------------------------------------------------

%Start of your document beginning with title page
\begin{document}


% --- Main Title Page ------------------------------------------------
\begin{titlepage}
\frontmatter

\begin{picture}(50,50)
\put(-70,40){\hbox{\includegraphics{images/logo.png}}}
\end{picture}

\vspace*{-5.8cm}

\begin{center}

\vspace{6.2cm}

\hspace*{-1.0cm} {\LARGE \textbf{Quality Control of Additive SLS Manufacturing\\}}
\vspace{0.2cm}
\hspace*{-1.0cm} using Convolutional Neural Networks\\

\vspace{2.0cm}

\hspace*{-1.0cm} { \textbf{Bachelor Thesis\\}}

\vspace{0.65cm}

\hspace*{-1.0cm} Submitted in partial fulfillment of the requirements for the degree of \\

\vspace{0.65cm}

\hspace*{-1.0cm} \textbf{Bachelor of Science in Engineering\\}

\vspace{0.65cm}

\hspace*{-1.0cm} to the University of Applied Sciences FH Campus Wien \\
\vspace{0.2cm}
\hspace*{-1.0cm} Bachelor Degree Program: Computer Science and Digital Communications \\

\vspace{1.6cm}

\hspace*{-1.0cm} \textbf{Author:} \\
\vspace{0.2cm}
\hspace*{-1.0cm} first name surname \\

\vspace{0.7cm}

\hspace*{-1.0cm} \textbf{Student identification number:}\\
\vspace{0.2cm}
\hspace*{-1.0cm} Number \\

\vspace{0.7cm}

\hspace*{-1.0cm} \textbf{Supervisor:} \\
\vspace{0.2cm}
\hspace*{-1.0cm} Title first name surname \\

\vspace{0.7cm}

% Reviewer if needed
%\hspace*{-1.0cm} \textbf{Reviewer: (optional)} \\
%\vspace{0.2cm}
%\hspace*{-1.0cm} Title first name surname \\


\vspace{1.0cm}

\hspace*{-1.0cm} \textbf{Date:} \\
\vspace{0.2cm}
\hspace*{-1.0cm} dd.mm.yyyy \\

\end{center}
\end{titlepage}

\newpage

\vspace*{16cm}
\setcounter{page}{1}

% --- Declaration of authorship ------------------------------------------
\hspace*{-0.7cm} \underline{Declaration of authorship:}\\\\
I declare that this Bachelor Thesis has been written by myself. I have not used any other than the listed sources, nor have I received any unauthorized help.\\\\
I hereby certify that I have not submitted this Bachelor Thesis in any form (to a reviewer for assessment) either in Austria or abroad.\\\\
Furthermore, I assure that the (printed and electronic) copies I have submitted are identical.
\\\\\\
Date: \hspace{6cm} Signature:\\

% --- English Abstract ----------------------------------------------------
\cleardoublepage
\chapter*{Abstract}
(E.g. ``This thesis investigates...'')
aaaa

% --- German Abstract ----------------------------------------------------
\cleardoublepage
\chapter*{Kurzfassung}
(Z.B. ``Diese Arbeit untersucht...'')


% --- Abbrevations ----------------------------------------------------
\chapter*{List of Abbreviations}
\vspace{0.65cm}

\begin{table*}[htbp]
		\begin{tabular}{ll}
			ARP & Address Resolution Protocol \\
			GPRS & General Packet Radio Service \\
			GSM  &  Global System for Mobile communication \\
			WLAN & Wireless Local Area Network \\
		\end{tabular}
\end{table*}

% --- Key terms ----------------------------------------------------
\newpage
\chapter*{Key Terms}
\vspace{0.65cm}

\begin{itemize}
	\setlength{\itemsep}{0pt}
	\item[] GSM
	\item[] Mobilfunk
	\item[] Zugriffsverfahren
\end{itemize}

% --- Table of contents autogenerated ------------------------------------
\newpage
\tableofcontents
\thispagestyle{empty}

% --- Begin of Thesis ----------------------------------------------------
\mainmatter
\chapter{Introduction}
\label{chap:intro}
Textkörper mit Bild

\begin{figure}[htbp]
	\centering
		\includegraphics[height=5cm]{images/buecher.png}
	\caption{Ein Stapel Bücher}
	\label{fig:buecher}
\end{figure}


Textkörper Fortsetzung mit Verweis auf den wundervollen Stapel Bücher in Abbildung \ref{fig:buecher}.


\section{Sub-chapter 1}
\label{sec:Unterkapitel1}

Textkörper mit Formel:

\begin{equation}
U(j\omega)=\int^{\infty}_{-\infty}{u(t) \cdot e^{-j\omega t}dt}
\label{form:form1}
\end{equation}

Textkörper Fortsetzung mit Verweis auf Formel \ref{form:form1}. Und nicht zu vergessen: es gibt auch noch eine tolle Abbildung in Kapitel \ref{chap:intro}, nämlich Abbildung \ref{fig:buecher}.

\subsection{Under sub-chapter 11}
\label{sec:UnterUnterkapitel11}

Textkörper mit direktem Zitat und Seitenanzahl:
``It would be very easy to show how technical or report writing differed from other writing'' \cite[p.~3]{young2002technical}.
\subsection{Under sub-chapter 12}
\label{sec:UnterUnterkapitel12}

Textkörper mit Referenzen:
Für weiterführende Informationen zum wissenschaftlichen Schreiben siehe "J. Schimel, Writing Science" \cite{schimel2012writing}. Es wird empfohlen den Sprachleitfaden der FH Campus Wien \cite{alker2006} zu berücksichtigen und die Checkliste für wissenschafltiches Schreiben \cite{petz2018} zu verwenden. Beide Leitfäden sind im FH Portal zu finden.

\chapter{Chapter 2}
\label{chap:back}

Textkörper mit noch einem Bild

\begin{figure}[htbp]
	\centering
		\includegraphics{images/birne}
	\caption{Eine Glühbirne}
	\label{fig:birne}
\end{figure}



\section{Sub-chapter 21}
\label{sec:Unterkapitel21}

Textkörper mit Tabelle.

\begin{table*}[htbp]
	\centering
		\begin{tabular}{|l|c|r|}
		\hline
		\rowcolor[gray]{0.9}
		Spalte 1 & Spalte 2 & Spalte 3 \\
		\hline
		Affen & Giraffen & Löwen \\
		Apfel & Birnen & Bananen \\
		Irgend & et & was \\
		\hline
		\end{tabular}
	\caption{Beispiel für eine Tabelle}
	\label{tab:BeispielFuerEineTabelle}
\end{table*}

Man beachte die Gegenüberstellung in Tabelle \ref{tab:BeispielFuerEineTabelle}.

%Online Tabellengenerator für Latex: https://www.tablesgenerator.com/

\section{Sub-chapter 23}
\label{sec:Unterkapitel23}

Aufzählungen:

Nummeriert:

\begin{enumerate}
	\item Punkt 1
	\item Punkt 2
\end{enumerate}

Mit Bullet Points:

\begin{itemize}
	\item Punkt 1
	\item Punkt 2
\end{itemize}

Mit Beschreibungen:

\begin{description}
	\item[Item 1] das ist der 1.Punkt
	\item[Item 2] und das der 2.
\end{description}


Auch Programmcodes können an entsprechender Stelle eingefügt werden, man beachte dazu auch Listing \ref{lst:conv}.

% see also http://mirror.easyname.at/ctan/macros/latex/contrib/listings/listings.pdf for options

\begin{lstlisting}[frame=lines, caption=Simple Listing, captionpos=b, label = lst:conv, language=C, showstringspaces=false]
#include <stdio.h>
int main()
{
	int i, n, t1 = 0, t2 = 1, nextTerm;

	printf("Enter the number of terms: ");
	scanf("%d", &n);

	printf("Fibonacci Series: ");

	for (i = 1; i <= n; ++i)
	{
		printf("%d, ", t1);
		nextTerm = t1 + t2;
		t1 = t2;
		t2 = nextTerm;
	}
	return 0;
}
\end{lstlisting}

Und zuguterletzt, Formeln mitten im Fliesstext, wie z.B. $a^2+b^2=c^2$, in einem Absatz.

\newpage
\chapter{Related Work}

\newpage
\chapter{Conclusion}

\newpage
\chapter{Future work}

\newpage

% --- Bibliography ------------------------------------------------------

%IEEE Citation [1]
\bibliographystyle{IEEEtran}
%for alphanumeric citation eg.: [ABC19]
%\bibliographystyle{alpha}

% List references I definitely want in the bibliography,
% regardless of whether or not I cite them in the thesis.

\newpage
\addcontentsline{toc}{chapter}{Bibliography}
\bibliography{thesisBib}

\newpage

% --- List of Figures ----------------------------------------------------

\addcontentsline{toc}{chapter}{List of Figures}
\listoffigures


% --- List of Tables -----------------------------------------------------

\newpage
\addcontentsline{toc}{chapter}{List of Tables}
\listoftables

% --- Appendix A -----------------------------------------------------

\backmatter
\appendix
\begin{appendices}
\chapter{Appendix}

(Hier können Schaltpläne, Programme usw. eingefügt werden.)

\clearpage
\end{appendices}

\end{document}
